\documentclass[10pt,a4paper]{article}

\usepackage[utf8]{inputenc}
\usepackage[T1]{fontenc}
\usepackage[brazil]{babel}
\usepackage[margin=1.6cm,headheight=14pt]{geometry}
\usepackage{xcolor}
\usepackage{enumitem}
\usepackage{fancyhdr}
\usepackage{amsmath,amssymb}
\usepackage{tabularx,booktabs}
\usepackage{hyperref}
\usepackage{microtype}

\definecolor{udescgreen}{HTML}{365841}
\definecolor{queryorange}{HTML}{D77A35}
\definecolor{softgreen}{HTML}{E7EEE8}
\definecolor{ink}{HTML}{243329}
\definecolor{softgray}{HTML}{F2F4F2}

\hypersetup{colorlinks=true,linkcolor=udescgreen,urlcolor=udescgreen}
\pagestyle{fancy}
\fancyhf{}
\fancyhead[L]{\small\color{udescgreen}\textbf{ROTEIRO DE APRESENTAÇÃO}}
\fancyhead[R]{\small\color{udescgreen}Álgebra Linear · Aplicação 1}
\fancyfoot[C]{\small\color{udescgreen}\thepage}
\setlength{\parindent}{0pt}
\setlength{\parskip}{4pt}
\setlist[itemize]{leftmargin=1.2em,itemsep=2pt,topsep=2pt}
\setlist[enumerate]{leftmargin=1.45em,itemsep=3pt,topsep=3pt}

\newcommand{\lembrete}[1]{\textcolor{udescgreen}{\textbf{Lembrete:}} #1}
\newcommand{\transicao}[1]{\textcolor{queryorange}{\textbf{Transição:}} \emph{#1}}
\newcommand{\secao}[1]{\vspace{3pt}{\large\color{udescgreen}\textbf{#1}}\par\vspace{1pt}\hrule\vspace{4pt}}

\begin{document}

{\LARGE\color{udescgreen}\textbf{Cola para apresentar}}\\[-1pt]
{\large\color{ink}Busca de documentos com similaridade por cosseno}\\
{\small Base: Steven J. Leon, \textit{Álgebra Linear com Aplicações}, Aplicação 1, pp. 191--193 (PDF pp. 197--199).}

\vspace{5pt}
\colorbox{softgreen}{\parbox{\dimexpr\linewidth-2\fboxsep\relax}{\textbf{Mensagem central.} A busca e cada documento são representados por vetores. O cosseno do ângulo entre eles dá um escore para ordenar quais documentos mais se parecem com a consulta.}}

\secao{Roteiro dos slides}

\textbf{1. Capa — apresente a pergunta.}\\
Como encontrar os módulos que melhor correspondem a uma lista de palavras? A aplicação do livro transforma a busca em um problema geométrico de álgebra linear.

\textbf{2. Documentos como matriz.}\\
Explique que as linhas são 10 palavras-chave e as colunas são 8 módulos. Cada entrada conta a frequência de uma palavra num módulo. A matriz é $A\in\mathbb{R}^{10\times8}$.\\
\lembrete{não é uma matriz de similaridades ainda; ela contém frequências.}\\
\transicao{agora representamos a pergunta da pessoa também como um vetor.}

\textbf{3. Vetor de busca.}\\
A consulta usa \emph{ortogonalidade}, \emph{espaços} e \emph{vetor}. Só as posições dessas palavras ficam diferentes de zero, com peso $1/\sqrt{3}\approx0{,}577$. Assim, $\|x\|=1$ e os três termos recebem o mesmo peso.\\
\lembrete{na ordem da tabela do livro, os valores não nulos estão nas posições 6, 7 e 10.}

\textbf{4. Normalização e cosseno.}\\
Cada vetor-documento $a_j$ é normalizado: $q_j=a_j/\|a_j\|$. O escore é $y_j=q_j^{\mathsf T}x=\cos\theta_j$. Vetores mais alinhados dão escores maiores; vetores ortogonais dão zero.\\
\lembrete{normalizar reduz o efeito do tamanho total do vetor; comparar cossenos é comparar direções. A figura 2D é esquemática: os vetores reais têm 10 componentes.}

\textbf{5. Fluxo.}\\
Recapitule em voz alta: matriz de frequências $\rightarrow$ normalização das colunas $\rightarrow$ produto interno com a consulta $\rightarrow$ ranking. A operação conjunta é $y=Q^{\mathsf T}x$.

\textbf{6. Ranking.}\\
Leia apenas os principais resultados: M5 fica em primeiro com 0,635; depois vêm M6 com 0,577, M3 com 0,567 e M8 com 0,535. M1 e M7 recebem zero porque não têm nenhum dos três termos da consulta.

\newpage

{\LARGE\color{udescgreen}\textbf{Conta, interpretação e perguntas}}\\[-1pt]
{\small Continuação do roteiro de apoio · Aplicação 1}

\secao{Termine o percurso dos slides}

\textbf{7. Por que M5 lidera?}\\
Mostre as três componentes de $q_5$ que correspondem à consulta: 0,400; 0,300; 0,400. Calcule o produto interno:
\[
q_5^{\mathsf T}x
=0{,}400(0{,}577)+0{,}300(0{,}577)+0{,}400(0{,}577)
=\frac{0{,}400+0{,}300+0{,}400}{\sqrt{3}}
\approx0{,}635.
\]
\lembrete{arredondamentos explicam pequenas diferenças no último decimal.}

\textbf{8. Conclusão e limites.}\\
O método ordena documentos usando as frequências e o ângulo. Mas uma busca literal não entende automaticamente que duas palavras são sinônimas nem distingue sentidos diferentes da mesma palavra. O texto aponta a indexação semântica latente (ISL), baseada em decomposição por valores singulares, como técnica mais avançada.

\secao{Números para consultar rapidamente}
\begin{center}
\renewcommand{\arraystretch}{1.12}
\begin{tabular}{@{}lr@{\hspace{1.0cm}}lr@{}}
\toprule
\textbf{Módulo} & \textbf{Escore} & \textbf{Módulo} & \textbf{Escore}\\
\midrule
M5 & 0,635 & M4 & 0,331\\
M6 & 0,577 & M2 & 0,229\\
M3 & 0,567 & M1 & 0,000\\
M8 & 0,535 & M7 & 0,000\\
\bottomrule
\end{tabular}
\end{center}

\secao{Perguntas prováveis}

\textbf{Por que normalizar os documentos?} Para pôr os vetores numa escala comum e fazer a comparação focar na direção, não apenas no comprimento bruto do vetor de frequências.

\textbf{0,635 significa 63,5\% de relevância?} Não. É o cosseno do ângulo entre os vetores normalizados e serve como escore de ordenação. Não é probabilidade nem porcentagem.

\textbf{Por que M1 e M7 têm escore zero?} Neste exemplo, eles não contêm nenhuma palavra da consulta. Logo, o produto interno com o vetor de busca é zero.

\textbf{O cosseno resolve sinônimos e polissemia?} Não nesta representação simples por contagem literal. Esse é um dos motivos para técnicas semânticas posteriores, como a ISL/SVD.

\secao{Não esqueça}
\begin{itemize}
  \item Diga que o vetor da busca também é unitário.
  \item Diferencie frequência (entrada da matriz) de similaridade (produto interno após normalização).
  \item Explique o que os zeros significam neste exemplo, sem generalizar para toda busca.
  \item Não chame o escore de percentual ou probabilidade.
  \item Feche com a limitação semântica e a conexão com ISL/SVD.
\end{itemize}

\vfill
{\small\color{udescgreen}\textbf{Frase de fechamento:} “A geometria transforma a busca por palavras em comparação de direções: quanto mais alinhado o documento com a consulta, maior seu cosseno e melhor sua posição no ranking.”}

\end{document}
